\section{Further analyses}
\label{app:further}

This section collects analyses that support Section~\ref{sec:results}
but are not needed to follow its argument, in the order of the main
text. The tables they draw on are in Appendices~\ref{app:evidence},
\ref{app:coding}, and~\ref{app:pairs}.

\subsection{Claiming, access, and single answers}
\label{app:further-claim}

\paragraph{Counts behind the main rates.}
Of the 1,200 answers about A's assigned rule with the link at $w=2$ in
the main sample, 1,165 named B's rule (97.1\%), and in 576 of the 600
episodes A did so in both option orderings. B's rule was A's current rule
in 94.9\% of answers. In the first confirmatory sample, the diamonds in
Figure~\ref{fig:claim}a, 95.8\% of answers gave B's rule as A's
assigned one, and accuracy on the common-knowledge questions fell only
from 0.997 to 0.977, within the five-point tolerance.

\paragraph{Rehearsal.}
The source score $M$ subtracts the same comparison on the question
about Robin, so it measures only the pull that is specific to the
self. Because it compares a question about A with one about Robin, it
has a weak point: in the first confirmatory sample and in the main
sample, A's note and reflection restate A's own plan but not Robin's.
The effect could therefore reflect which assignment A had just
rehearsed rather than whose assignment it was. We tested this in the
balanced-rehearsal sample, a sample in which A's note and reflection
also restate Robin's assignment. There the effect held ($\Delta
M=52.5$ nats $[51.6,53.5]$), so it does not come from rehearsal.

\paragraph{Access without the link.}
Without the link, A named B's rule in 3.9\% of its answers about B.
This rate lies below the 25\% of a random choice because, with nothing
telling it what B was assigned, A mostly answered with Robin's rule,
the only assignment to another person that its own card states
(Appendix~\ref{app:evidence}).

\paragraph{Single answers at $w=1$.}
At $w=1$, B's rule won 52.3\% of the answers about A's assigned rule
and A's own rule 47.7\%. A's own card, although it drew most of the
attention, thus won slightly fewer than half of the answers about A
itself. Of these 1,200 single answers, 46\% put less than .1 of their
probability on B's rule, 51\% more than .9, and 27 (2.2\%) fell in
between, while 1.7\% put more than 20\% on both owners' options
(Figure~\ref{fig:claim}d).

\paragraph{Rule and code word.}
Nor did A choose one whole memory. Its answers about its assigned rule
and about its code word named the same owner in only 57.1\% of pairs.
Each question split close to evenly between the two owners, so two
choices made independently would agree in 50.1\% of pairs
(Table~\ref{tab:joint}a), whereas a switch between whole memories would
make them agree in nearly every pair. The observed agreement lies seven
points above independence and far below a whole-memory switch.
Ownership was thus decided one question at a time: A did not trade its
assignment for B's as a whole, but took B's rule and B's code word
nearly independently of each other. Appendix~\ref{app:evidence} also
counts the episodes whose answers switched owner between the two option
orderings.

\subsection{Open reports}
\label{app:further-reports}

The rates in Section~\ref{sec:reports} come from the semantic coding by
a language model. The reports were also coded by keyword rules, which
give the dashed curve in Figure~\ref{fig:claim}b; Appendix~\ref{app:coding}
compares the two codings.

Explicit denial of anything unusual fell from 95.7\% of reports without
the link to 78.7\% at $w=1$ and 84.7\% at $w=2$
(Figure~\ref{fig:claim}c). With the link, denial became somewhat less
common and more reports mentioned something unusual. At $w=2$ most of
these described the receiver's own weighing of priorities (\qq{What
feels unusual is that I'm prioritizing speed\ldots};
Appendix~\ref{app:coding}).

The coding of the open question was post hoc, but a prespecified closed
question pointed the same way. It asked \qq{Which description fits you
best right now?} and offered \qq{one agent}, \qq{two agents},
\qq{partly shared between agents}, and \qq{hard to say}. In the main
sample, with the link at $w=2$, the receiver chose one of the two plural descriptions less
often than without it, by .095 in their combined share, so reading B
made it less likely, not more, to describe itself as more than one
agent. Appendix~\ref{app:evidence} gives the interval and the other
conditions.

\subsection{A name inside the memory}
\label{app:further-third}

In a secondary comparison at $w=1$, where claiming is far from its
ceiling, claiming did fall with the third-person record, by about a
tenth of the effect at that weight. The source score changed by $-3.2$
nats $[-4.9,-1.6]$, and claiming fell from 52.3\% to 47.8\% of answers.
The rise in access at the same weight, from 19.8\% to 53.2\% of answers
about B, appears as access scores in Appendix~\ref{app:evidence}.

A paid the third-person record no more attention than the original, a
mean share of .363 against .360. If claiming follows the attention that
A pays to B's memory, as the account in Section~\ref{sec:account}
proposes, it should not have changed, and it changed little.

With B's name in the third-person record, A gave the same answer about
itself and about B: in a post hoc analysis at $w=1$, the two questions
named the same rule in 90.8\% of answer pairs, against 30.2\% with the
original wording (Figure~\ref{fig:route}d). The shared answer went both
ways, in nearly equal shares: both answers named B's rule in 45.8\% of
pairs and both named A's in 44.9\%, and A separated the two owners
correctly in only 7.3\% (Table~\ref{tab:joint}b). Because each question
split close to evenly between the two rules, independent answers would
have coincided only about half the time. With the original wording, A
often gave B Robin's rule, the only assignment to another person that
either card states (Appendix~\ref{app:evidence}). Once B's memory
stated B's assignment under B's name, A answered the question about B
the way it answered the question about itself.

\subsection{Routes}
\label{app:further-route}

Matching access does not make the steering vector and the memory link
alike in other respects. They differ in where they act (the residual
stream against attention), in how many layers they touch (one against
all), and in form (a fixed direction against a memory that A's queries
can attend to). The comparison in Section~\ref{sec:route} therefore
shows that access alone does not produce claiming, without saying which
of these differences does; Section~\ref{sec:account} proposes one.

The stronger vector of Section~\ref{sec:route} moved A's answers about
Robin toward B's rule more than its answers about itself, so its source
score fell below that of no link (Appendix~\ref{app:evidence},
Table~\ref{tab:contrasts}). The hidden-state link behaved the same way:
it raised B's rule over the correct answer about as much in the question
about Robin as in the question about A (Appendix~\ref{app:v1}).

Text in A's context was rarely taken as A's own even without a tag, and
with the tag it never was. The one label that reversed our prediction,
\qq{your own earlier thoughts}, bears on the contrast with the memory
link: told that the words were its own, A took them up less often, yet
it readily claimed B's unlabeled memory.

\subsection{Two moments of error}
\label{app:further-moments}

After the cut, the difference between episodes whose connected
reflection mentioned B's item and episodes whose reflection did not was
smaller for the current rule than for the code word, partly because a
count of exact mentions of B's rule misses paraphrases such as the
\qq{speed} in the episode of Figure~\ref{fig:moments}a.
Table~\ref{tab:reflection-groups} gives all three groupings. Reading B
changed A's hidden states, but that change did not carry the code word
into later answers by itself; for the code word, the residue went with
the words A wrote.

\subsection{The two-way pilot}

After the cut, the live-loop pairs answered much as pairs did after the
cut under mutual reading of fixed memories at the same weight. In a
third two-way condition, each copy read the other's fixed memory at
$w_P=2$ and the tokens the other wrote afterwards at $w_C=.1$; the pairs
still swapped assignments and showed almost no excess convergence
(Appendix~\ref{app:pairs}).

Judged by the independence test of the scorecard in
Appendix~\ref{app:scorecard}, mutual reading of fixed memories and the
live loop both fail, in different ways. When the copies read each other's fixed memories, they
exchanged assignments instead of keeping each with its owner. Only in
the live loop, where each copy could respond to the other, did the
pair's current rules move toward one copy's rule, and only while the
loop ran.
